% Fragmento público generado desde propietarios íntegros.
% Residencia: 52.9.1.
% Fuente: ../incoming/radion_monografia_20260827/manuscrito/sections/10_minkowski_hodge_lorentz.tex:56-100
\subsubsection{Dualité de Hodge et opérateur de quart de tour}

La signature \((-+++ )\) et une orientation étant fixées, l'étoile de Hodge sur
les deux-formes satisfait
\begin{equation}
\star:\Lambda^2Q_4^*\to\Lambda^2Q_4^*,
\qquad \star^2=-I.
\label{eq:hodge-star}
\end{equation}
De cette manière, le quart de tour qui, dans le plan de phase, s'exprimait au moyen de
\(J^2=-I\) réapparaît dans l'espace des champs comme structure complexe sur
\(\Lambda^2\). Les domaines sont différents, mais la généalogie d'orientation
est conservée par le foncteur.

Soit \(A_{\mathrm{em}}\) un potentiel et
\begin{equation}
F=\diff A_{\mathrm{em}}.
\label{eq:F-dA}
\end{equation}
La décomposition de \(F\) par rapport à un observateur temporel produit des champs
électrique et magnétique ; \(\star F\) échange leurs rôles avec le signe fixé
par l'orientation et la signature. La polarité constitutive de
\eqref{eq:EM-conjugation} acquiert ainsi une réalisation différentielle.

\paragraph{Chaîne de transduction électromagnétique et causale}

Le transport qui doit rester visible est
\begin{equation}
\Gamma_9
\longrightarrow\text{hélice de mémoire}
\longrightarrow(Q_4,B_{\mathrm{caus}},\star)
\longrightarrow A_{\mathrm{em}}
\longrightarrow F=\diff A_{\mathrm{em}}
\longrightarrow
m\nabla_u u=q(\iota_uF)^\sharp.
\label{eq:lorentz-chain}
\end{equation}
Chaque flèche ajoute de la structure : l'holonomie conserve le retour ; le quotient fixe
la causalité ; Hodge fixe la dualité ; le potentiel fixe le champ ; le couplage
avec la charge produit la dynamique.

L'équation finale est la force de Lorentz sous forme géométrique. L'hélice TPK
est son antécédent de mémoire, non une orbite physique déjà mue par un champ. La
dynamique n'apparaît qu'à l'introduction de \(F\), \(q\) et de la connexion \(\nabla\).

